%%% ============================================================
%%% اینفوگرافیک دو صفحه‌ای — طرح مجلس مهستان
%%% فرمت: A4 عمودی — قابل چاپ و اشتراک‌گذاری
%%% ============================================================

\documentclass[12pt, a4paper]{article}
\usepackage[margin=1cm]{geometry}
\usepackage{tikz}
\usepackage{xcolor}
\usepackage{graphicx}

\usetikzlibrary{
	shapes.geometric,
	shapes.misc,
	arrows.meta,
	positioning,
	calc,
	backgrounds,
	fit,
	decorations.pathmorphing,
	shadows.blur
}

\usepackage{fontspec}
\usepackage{xepersian}

\settextfont[Scale=1.1]{Vazirmatn}
\setdigitfont[Scale=1.1]{Vazirmatn}
\setlatintextfont{Times New Roman}

\usepackage{hyperref}
\hypersetup{pdfencoding=unicode}

%% رنگ‌ها
\definecolor{DeepBlue}{HTML}{1A237E}
\definecolor{PrimaryBlue}{HTML}{1565C0}
\definecolor{LightBlue}{HTML}{BBDEFB}
\definecolor{PrimaryGreen}{HTML}{2E7D32}
\definecolor{LightGreen}{HTML}{C8E6C9}
\definecolor{PrimaryRed}{HTML}{C62828}
\definecolor{LightRed}{HTML}{FFCDD2}
\definecolor{PrimaryPurple}{HTML}{6A1B9A}
\definecolor{LightPurple}{HTML}{F3E5F5}
\definecolor{GoldYellow}{HTML}{F9A825}
\definecolor{LightYellow}{HTML}{FFF9C4}
\definecolor{PrimaryOrange}{HTML}{E65100}
\definecolor{LightOrange}{HTML}{FFF3E0}
\definecolor{TealColor}{HTML}{00695C}
\definecolor{LightTeal}{HTML}{E0F2F1}
\definecolor{DarkGray}{HTML}{424242}
\definecolor{LightGray}{HTML}{F5F5F5}

\pagestyle{empty}

\begin{document}
	
	%% ══════════════════════════════════════════════════════
	%% صفحه‌ی اول
	%% ══════════════════════════════════════════════════════
		\centering
		\begin{tikzpicture}[remember picture, overlay]
		
		%% پس‌زمینه
		\fill[LightGray] (current page.south west) rectangle (current page.north east);
		
		%% هدر
		\fill[DeepBlue]
		([yshift=-3cm]current page.north west) rectangle (current page.north east);
		\fill[GoldYellow]
		([yshift=-3cm]current page.north west) rectangle
		([yshift=-3.15cm]current page.north east);
		
		%% عنوان اصلی
		\node[anchor=center, text=white, font=\fontsize{28}{36}\selectfont\bfseries]
		at ([yshift=-1.2cm]current page.north)
		{\rl{طرح مجلس مهستان}};
		\node[anchor=center, text=GoldYellow, font=\large]
		at ([yshift=-2.2cm]current page.north)
		{\rl{چارچوب نهادی گذار دموکراتیک ایران}};
		
		%% ═══ بخش ۱: چرا مهستان؟ ═══

		\node[anchor=north west, font=\Large\bfseries\color{DeepBlue}]
		at ([xshift=80mm, yshift=-35mm]current page.north west)
		{\rl{چرا این طرح؟}};
		
		\centering
		\draw[DeepBlue, line width=1.5pt]
		([xshift=7.75cm,  yshift=-4.4cm]current page.north west) --
		([xshift=12.75cm,  yshift=-4.4cm]current page.north west);
		
		%% باکس‌های مشکل و راه‌حل
		\node[rectangle, rounded corners=8pt, fill=LightRed, draw=PrimaryRed,
		text width=8.5cm, align=right, inner sep=10pt, font=\small]
		at ([xshift=5.5cm, yshift=-6cm]current page.north west)
		{%
			\textbf{\color{PrimaryRed}\rl{مشکل: خلأ نهادی}}\\[4pt]
			\rl{تجربه‌ی لیبی، مصر و عراق نشان داد:}\\
			\rl{بدون نهاد واحد و مشروع = هرج‌ومرج}%
		};
		
		\node[rectangle, rounded corners=8pt, fill=LightGreen, draw=PrimaryGreen,
		text width=8.5cm, align=right, inner sep=10pt, font=\small]
		at ([xshift=15.5cm, yshift=-6cm]current page.north west)
		{%
			\textbf{\color{PrimaryGreen}\rl{راه‌حل: مجلس مهستان}}\\[4pt]
			\rl{۳۰۰ نماینده‌ی منتخب مستقیم مردم}\\
			\rl{نهاد عالی حاکمیت دوران گذار}%
		};
		
		\draw[-Stealth, line width=2pt, PrimaryGreen]
		([xshift=10cm, yshift=-6cm]current page.north west) --
		([xshift=11cm, yshift=-6cm]current page.north west);
		
		%% ═══ بخش ۲: ترکیب مجلس ═══
		\node[anchor=north west, font=\Large\bfseries\color{DeepBlue}]
		at ([xshift=7cm, yshift=-8cm]current page.north west)
		{\rl{ترکیب مجلس مهستان}};
		
		\draw[DeepBlue, line width=1.5pt]
		([xshift=7.75cm, yshift=-8.9cm]current page.north west) --
		([xshift=12.75cm, yshift=-8.9cm]current page.north west);
		
		%% نمودار دایره‌ای
		\begin{scope}[shift={([xshift=5cm, yshift=-12.5cm]current page.north west)}]
			\fill[DeepBlue] (0,0) circle (1.8cm);
			\node[text=white, font=\Large\bfseries, align=center] at (0,0)
			{\rl{۳۰۰}\\[-2pt]{\small\rl{نماینده}}};
			\fill[PrimaryBlue]  (0:2.2cm)   arc (0:288:2.2cm)   -- (288:1.9cm) arc (288:0:1.9cm)   -- cycle;
			\fill[GoldYellow]   (288:2.2cm) arc (288:336:2.2cm) -- (336:1.9cm) arc (336:288:1.9cm) -- cycle;
			\fill[PrimaryPurple](336:2.2cm) arc (336:360:2.2cm) -- (360:1.9cm) arc (360:336:1.9cm) -- cycle;
		\end{scope}
		
		%% لیبل‌های ترکیب
		\node[rectangle, rounded corners=5pt, fill=LightBlue, draw=PrimaryBlue,
		text width=4cm, align=center, inner sep=6pt, font=\small]
		at ([xshift=11cm, yshift=-10.5cm]current page.north west)
		{\textbf{\color{PrimaryBlue}\rl{۲۴۰ کرسی استانی}}\\[2pt]
			\rl{سیستم نسبی-فهرستی}};
		
		\node[rectangle, rounded corners=5pt, fill=LightYellow, draw=GoldYellow,
		text width=4cm, align=center, inner sep=6pt, font=\small]
		at ([xshift=11cm, yshift=-12.5cm]current page.north west)
		{\textbf{\color{GoldYellow!80!black}\rl{۴۰ کرسی ملی}}\\[2pt]
			\rl{فهرست سراسری}};
		
		\node[rectangle, rounded corners=5pt, fill=LightPurple, draw=PrimaryPurple,
		text width=4cm, align=center, inner sep=6pt, font=\small]
		at ([xshift=11cm, yshift=-14.5cm]current page.north west)
		{\textbf{\color{PrimaryPurple}\rl{۲۰ کرسی اقلیت‌ها}}\\[2pt]
			\rl{تضمینی}};
		
		\node[rectangle, rounded corners=5pt, fill=LightRed, draw=PrimaryRed,
		text width=4cm, align=center, inner sep=6pt, font=\small]
		at ([xshift=16.5cm, yshift=-12.5cm]current page.north west)
		{\textbf{\color{PrimaryRed}\rl{حداقل ۳۰٪}}\\[2pt]
			\rl{تنوع جنسیتی}\\
			\rl{در فهرست‌ها}};
		
		%% ═══ بخش ۳: ۶ اصل راهنما ═══
		%% عنوان — وسط صفحه (xshift≈10.25cm)
		\node[anchor=north, font=\Large\bfseries\color{DeepBlue}]
		at ([xshift=10.25cm, yshift=-16cm]current page.north west)
		{\rl{۶ اصل راهنما}};
		
		\draw[DeepBlue, line width=1.5pt]
		([xshift=7.75cm, yshift=-16.9cm]current page.north west) --
		([xshift=12.75cm, yshift=-16.9cm]current page.north west);
		
		%% ردیف اول — سه دایره، مرکز‌چین حول xshift=10.25cm
		%% فاصله‌ی مرکزها: 5cm — بازه: 5cm تا 10.25 تا 15.5cm
		\foreach \i/\title/\col in {
			1/{\rl{دموکراسی + تخصص}}/PrimaryBlue,
			2/{\rl{نمایندگی نسبی}}/PrimaryGreen,
			3/{\rl{دو قانون اساسی}}/PrimaryPurple%
		}{
			\node[circle, fill=\col, minimum size=1.2cm,
			font=\large\bfseries\color{white}]
			at ([xshift={5.25cm+(\i-1)*5cm}, yshift=-18.5cm]current page.north west)
			{\rl{\i}};
			\node[anchor=north, font=\small\bfseries\color{\col},
			text width=4cm, align=center]
			at ([xshift={5.25cm+(\i-1)*5cm}, yshift=-19.3cm]current page.north west)
			{\title};
		}
		
		%% ردیف دوم — سه دایره، همان محور مرکزی
		\foreach \i/\title/\col in {
			4/{\rl{اصلاح نه تخریب}}/PrimaryOrange,
			5/{\rl{شفافیت حداکثری}}/TealColor,
			6/{\rl{فراگیری و حقوق بنیادین}}/PrimaryRed%
		}{
			\node[circle, fill=\col, minimum size=1.2cm,
			font=\large\bfseries\color{white}]
			at ([xshift={5.25cm+(\i-4)*5cm}, yshift=-21cm]current page.north west)
			{\rl{\i}};
			\node[anchor=north, font=\small\bfseries\color{\col},
			text width=4cm, align=center]
			at ([xshift={5.25cm+(\i-4)*5cm}, yshift=-21.8cm]current page.north west)
			{\title};
		}
		
		%% ═══ بخش ۴: خط زمانی ═══
		\node[anchor=north west, font=\Large\bfseries\color{DeepBlue}]
		at ([xshift=82mm, yshift=-23cm]current page.north west)
		{\rl{خط زمانی گذار}};
		
		\draw[DeepBlue, line width=1.5pt]
		([xshift=7.75cm,   yshift=-23.9cm]current page.north west) --
		([xshift=12.75cm, yshift=-23.9cm]current page.north west);
		
		\draw[DarkGray, line width=2pt, -{Stealth[length=4mm]}]
		([xshift=2cm,    yshift=-25.5cm]current page.north west) --
		([xshift=18.5cm, yshift=-25.5cm]current page.north west);
		
		\foreach \x/\phase/\col/\label in {
			2.5/۱/PrimaryOrange/{\rl{انتخابات}},
			5.5/۲/PrimaryRed/{\rl{اقدامات فوری}},
			8.5/۳/PrimaryPurple/{\rl{ق.ا. موقت}},
			11.5/۴/PrimaryBlue/{\rl{نهادسازی}},
			14.5/۵/TealColor/{\rl{ق.ا. نهایی}},
			17.5/۶/PrimaryGreen/{\rl{انتقال قدرت}}%
		}{
			\fill[\col]
			([xshift=\x cm, yshift=-25.5cm]current page.north west) circle (6pt);
			\node[anchor=south, font=\scriptsize\bfseries\color{\col}]
			at ([xshift=\x cm, yshift=-25.3cm]current page.north west)
			{\rl{فاز \phase}};
			\node[anchor=north, font=\tiny\color{DarkGray}]
			at ([xshift=\x cm, yshift=-25.7cm]current page.north west)
			{\label};
		}
		
		\node[font=\scriptsize\color{DarkGray}]
		at ([xshift=2.5cm,  yshift=-26.3cm]current page.north west) {\rl{هفته ۵}};
		\node[font=\scriptsize\color{DarkGray}]
		at ([xshift=8.5cm,  yshift=-26.3cm]current page.north west) {\rl{هفته ۱۰}};
		\node[font=\scriptsize\color{DarkGray}]
		at ([xshift=14.5cm, yshift=-26.3cm]current page.north west) {\rl{ماه ۱۸}};
		\node[font=\scriptsize\color{DarkGray}]
		at ([xshift=17.5cm, yshift=-26.3cm]current page.north west) {\rl{ماه ۲۴}};
		
		%% فوتر
		\fill[DeepBlue]
		([yshift=1.2cm]current page.south west) rectangle (current page.south east);
		\fill[GoldYellow]
		([yshift=1.2cm]current page.south west) rectangle
		([yshift=1.35cm]current page.south east);
		
		\node[anchor=center, text=white, font=\normalsize]
		at ([yshift=0.6cm]current page.south)
		{\rl{صفحه‌ی ۱ از ۲} \quad | \quad \rl{نسخه‌ی ۲.۰ — اسفند ۱۴۰۴}};
		
	\end{tikzpicture}
	
	\newpage
	
	%% ══════════════════════════════════════════════════════
	%% صفحه‌ی دوم
	%% ══════════════════════════════════════════════════════
	\begin{tikzpicture}[remember picture, overlay]
		
		%% پس‌زمینه
		\fill[LightGray] (current page.south west) rectangle (current page.north east);
		
		%% هدر
		\fill[DeepBlue]
		([yshift=-2.5cm]current page.north west) rectangle (current page.north east);
		\fill[GoldYellow]
		([yshift=-2.5cm]current page.north west) rectangle
		([yshift=-2.65cm]current page.north east);
		
		\node[anchor=center, text=white, font=\fontsize{22}{30}\selectfont\bfseries]
		at ([yshift=-0.75cm]current page.north)
		{\rl{عدالت انتقالی و اقدامات فوری}};
		
		%% ═══ بخش ۱: ۵ ستون عدالت انتقالی ═══
		%% عنوان — وسط‌چین، 50mm بالاتر از قبل (yshift از -3.5cm به -0.5cm)
		\node[anchor=north, font=\Large\bfseries\color{DeepBlue}]
		at ([xshift=9.75cm, yshift=-30mm]current page.north west)
		{\rl{۵ ستون عدالت انتقالی}};
		
		\draw[DeepBlue, line width=1.5pt]
		([xshift=6.25cm, yshift=-4cm]current page.north west) --
		([xshift=14.25cm, yshift=-4cm]current page.north west);
		
		%% ستون‌ها — مرکز‌چین (xshift مرکز ۵ باکس با فاصله ۳.۵cm = ۱.۷۵ تا ۱۵.۷۵، مرکز ۸.۷۵cm)
		%% برای وسط‌چین حول ۱۰.۲۵cm: offset = 10.25 - 8.75 = 1.5cm → شروع از 3.25cm
		\foreach \i/\title/\col/\items in {
			1/{\rl{حقیقت‌یابی}}/PrimaryBlue/{%
				\rl{کمیسیون حقیقت}\\
				\rl{مستندسازی جنایات}\\
				\rl{ثبت شهادت‌ها}},
			2/{\rl{عدالت کیفری}}/PrimaryRed/{%
				\rl{محاکمات عادلانه}\\
				\rl{تفکیک مسئولیت‌ها}\\
				\rl{عفو مشروط}},
			3/{\rl{جبران خسارت}}/PrimaryOrange/{%
				\rl{غرامت مالی}\\
				\rl{اعاده‌ی حیثیت}\\
				\rl{خدمات درمانی}},
			4/{\rl{آشتی ملی}}/GoldYellow/{%
				\rl{گفتگوهای ملی}\\
				\rl{موزه‌ها و یادبود}\\
				\rl{اصلاح تاریخ‌نگاری}},
			5/{\rl{عدم تکرار}}/PrimaryGreen/{%
				\rl{اصلاح نهادی}\\
				\rl{آموزش حقوق بشر}\\
				\rl{نهادهای نظارتی}}%
		}{
			\node[rectangle, rounded corners=6pt,
			fill=\col!15, draw=\col, line width=1pt,
			text width=3.2cm, align=center,
			minimum height=3.5cm, inner sep=6pt]
			at ([xshift={3.25cm+(\i-1)*3.5cm}, yshift=-6cm]current page.north west)
			{%
				\textbf{\color{\col}\title}\\[4pt]
				\rule{2cm}{0.3pt}\\[3pt]
				{\scriptsize\items}%
			};
		}
		
		%% ═══ بخش ۲: ۱۰ اقدام فوری ═══
		%% عنوان — وسط‌چین، 50mm بالاتر (yshift از -9.5cm به -7cm)
		\node[anchor=north, font=\Large\bfseries\color{DeepBlue}]
		at ([xshift=9.75cm, yshift=-8cm]current page.north west)
		{\rl{۱۰ اقدام فوری — هفته‌ی اول}};
		
		\draw[DeepBlue, line width=1.5pt]
		([xshift=5.75cm, yshift=-8.95cm]current page.north west) --
		([xshift=14.75cm, yshift=-8.95cm]current page.north west);
		
		%% ستون چپ (۱–۵) — مرکز تقریبی xshift=5cm
		\foreach \i/\action in {
			1/{\rl{آزادی زندانیان سیاسی}},
			2/{\rl{لغو حجاب اجباری}},
			3/{\rl{رفع فیلترینگ اینترنت}},
			4/{\rl{تعلیق مجازات اعدام}},
			5/{\rl{آزادی فعالیت احزاب}}%
		}{
			\node[circle, fill=PrimaryGreen, minimum size=0.7cm,
			font=\small\bfseries\color{white}]
			at ([xshift=2cm, yshift={-9.6cm-(\i-1)*0.9cm}]current page.north west)
			{\rl{\i}};
			\node[anchor=west, font=\small]
			at ([xshift=2.6cm, yshift={-9.6cm-(\i-1)*0.9cm}]current page.north west)
			{\action};
		}
		
		%% ستون راست (۶–۱۰) — مرکز تقریبی xshift=13cm
		\foreach \i/\action in {
			6/{\rl{ممنوعیت بازداشت بدون حکم}},
			7/{\rl{لغو محاکم انقلاب}},
			8/{\rl{حق بازگشت تبعیدیان}},
			9/{\rl{مسدودسازی حساب‌های مشکوک}},
			10/{\rl{حفاظت از اسناد امنیتی}}%
		}{
			\node[circle, fill=PrimaryGreen, minimum size=0.7cm,
			font=\small\bfseries\color{white}]
			at ([xshift=11cm, yshift={-9.6cm-(\i-6)*0.9cm}]current page.north west)
			{\rl{\i}};
			\node[anchor=west, font=\small]
			at ([xshift=11.6cm, yshift={-9.6cm-(\i-6)*0.9cm}]current page.north west)
			{\action};
		}
		
		%% ═══ بخش ۳: ۱۴ اصل غیرقابل‌نقض ═══
		%% عنوان — وسط‌چین، 50mm بالاتر (yshift از -16cm به -14.5cm)
		\node[anchor=north, font=\Large\bfseries\color{DeepBlue}]
		at ([xshift=10.25cm, yshift=-14cm]current page.north west)
		{\rl{۱۴ اصل غیرقابل‌نقض}};
		
		\draw[DeepBlue, line width=1.5pt]
		([xshift=6.25cm, yshift=-14.95cm]current page.north west) --
		([xshift=14.25cm, yshift=-14.95cm]current page.north west);
		
		\node[rectangle, rounded corners=8pt,
		fill=LightPurple, draw=PrimaryPurple, line width=1.5pt,
		text width=17.5cm, align=center, inner sep=12pt]
		at ([xshift=10.25cm, yshift=-17.5cm]current page.north west)
		{%
			{\small%
				\begin{tabular}{r@{\hspace{0.8cm}}r@{\hspace{0.8cm}}r}
					\rl{۱. حاکمیت مردم}     & \rl{۲. جدایی دین از دولت} & \rl{۳. برابری}            \\[3pt]
					\rl{۴. آزادی‌های بنیادین}& \rl{۵. ممنوعیت شکنجه}     & \rl{۶. دادرسی عادلانه}    \\[3pt]
					\rl{۷. حقوق اقلیت‌ها}   & \rl{۸. حقوق زنان}         & \rl{۹. استقلال قضایی}     \\[3pt]
					\rl{۱۰. تمامیت ارضی}    & \rl{۱۱. حاکمیت قانون}     & \rl{۱۲. عدالت نه انتقام} \\[3pt]
					\rl{۱۳. آزادی اطلاعات}  & \rl{۱۴. تعهد به صلح}      &                           \\
				\end{tabular}%
			}%
		};
		
		%% ═══ بخش ۴: فراخوان ═══
		\node[rectangle, rounded corners=10pt,
		fill=LightGreen, draw=PrimaryGreen, line width=2pt,
		text width=17.5cm, align=center, inner sep=15pt]
		at ([xshift=10.25cm, yshift=-22cm]current page.north west)
		{%
			{\Large\bfseries\color{PrimaryGreen}\rl{فراخوان به همه‌ی ایرانیان}}\\[8pt]
			{\normalsize
				\textbf{\rl{هر جریان:}} \rl{قانون اساسی موقت خود را بنویسید} \quad
				\textbf{\rl{هر شهروند:}} \rl{مطالعه و مشارکت کنید}\\[4pt]
				\textbf{\rl{هر متخصص:}} \rl{دانش خود را در خدمت گذار بگذارید} \quad
				\textbf{\rl{هر فعال:}} \rl{مستندسازی را ادامه دهید}%
			}%
		};
		
		%% شعار نهایی
		\node[anchor=center, font=\fontsize{20}{28}\selectfont\bfseries\color{DeepBlue}]
		at ([yshift=-25cm]current page.north)
		{\rl{«ایران برای همه‌ی ایرانیان»}};
		
		%% فوتر
		\fill[DeepBlue]
		([yshift=1.2cm]current page.south west) rectangle (current page.south east);
		\fill[GoldYellow]
		([yshift=1.2cm]current page.south west) rectangle
		([yshift=1.35cm]current page.south east);
		
		\node[anchor=center, text=white, font=\normalsize]
		at ([yshift=0.6cm]current page.south)
		{\rl{صفحه‌ی ۲ از ۲} \quad | \quad
			\rl{تصمیم نهایی فقط با مردم ایران است}};
		
	\end{tikzpicture}
	
\end{document}